\section{A collar with round leaves}\label{sec:collar}

Our next task is to replace the large leaves by round spheres while
retaining the original data on their inner side. The metric change is
small, but smallness alone does not preserve a saturated dominant energy
condition. Instead, we arrange that the change in the constraint vector
is itself future causal.

The construction has two stages on adjacent annuli. On the first we
start a positive addition to $F$ and a radial addition to $L$. On the
second, with these additions fully active, we interpolate the remaining
fields to their round values. The first addition controls the small
null component $C-Q$; the second controls $C+Q$. Their product will
dominate the square of the angular momentum change. This is why no
strictness assumption on the original DEC is needed.

Throughout this section a subscript $0$ denotes the quantities of the
original data in the balanced chart. Constants denoted $C_*$ or $C_j$
may depend on these data and on fixed cutoff profiles, but not on the
large parameter $R$. They are distinct from the constraint density $C$
and its original value $C_0$.

\begin{lemma}[Round-leaf collar]\label{lem:collar}
Let $(g,K)$ be smooth data in the balanced class \eqref{eq:cks},
satisfying the dominant energy condition. For all sufficiently large
$R$, there are smooth data $(G,k)$ on the original manifold up to the
sphere $r=3R$ with the following properties.
They equal $(g,K)$ on the entire inner side of $r=R$, satisfy the dominant
energy condition, and on $R\leq r\leq3R$ satisfy
\[
 |G-h|_h\leq C_* r^{-3}.
\]
On their entire domain they also satisfy
$G\geq(1-C_*R^{-3})g$, after enlarging $C_*$ if necessary.
The constants and the threshold for $R$ depend on the given data and
fixed cutoff profiles, not on $R$.
At $r=3R$ their full jets agree with the prescription
\[
\begin{gathered}
 \gamma=r^2\sigma,\qquad s=\eta=\tau=0,\qquad
 v=r,\qquad F=f_*+c_R(r),\\
 L=1+r^{-2},\qquad U=u=\sqrt{1+r^2-2F/r}.
\end{gathered}
\]
Here $f_*=M/4$ is the balanced mass function, and $c_R$ is smooth,
nonnegative, and satisfies
\begin{equation}\label{eq:mass-padding}
 c_R(r)=\int_R^r w(\rho)\rho^{-3/2}\,\mathrm d\rho,
 \qquad c_R'(r)=w(r)r^{-3/2},\qquad
 0\leq c_R(r)\leq2R^{-1/2}\quad(r\geq R),
\end{equation}
where $w$ is nondecreasing, vanishes for $r\leq R$, and equals one for
$r\geq2R$.  In particular $c_R$ extends by zero across $r=R$.
\end{lemma}

\begin{proof}
Choose smooth nondecreasing cutoff functions $w,e$ from $0$ to $1$, each
a fixed profile of $r/R$, so that
\[
\begin{array}{c|cc}
 &\text{identically zero}&\text{identically one}\\ \hline
 w&r\leq R&r\geq2R\\
 e&r\leq2R&r\geq3R.
\end{array}
\]
Use \eqref{eq:mass-padding} to define $c_R$ for all $r\geq R$.  On the
collar $R\leq r\leq3R$ prescribe
\begin{equation}\label{eq:collar-data}
\begin{aligned}
 \gamma&=(1-e)\gamma_0+e r^2\sigma,&
 s&=(1-e)s_0,\\
 v&=(1-e)v_0+er,&
 F&=(1-e)F_0+ef_*+c_R,\\
 L_b&=(1-e)L_0+e,&
 L&=L_b+\frac{w}{r^2},\\
 \eta&=(1-e)\eta_0,&
 \tau&=(1-e)\left(\tau_0-
       \frac12(\tr_\gamma\tau_0)\gamma\right).
\end{aligned}
\end{equation}
The last expression is trace-free with respect to the new sphere metric.
Define $d$ from $\gamma,s$ as in \eqref{eq:mass-variables}, and then set
\begin{equation}\label{eq:collar-reconstruction}
 u=\sqrt{1+v^2-2F/r},\qquad U=u/d,\qquad
 k_T=(v/r)\gamma+\tau.
\end{equation}
Together with $k(N,N)=L$ and $k(N,\cdot)|_{T\Sph}=\eta$, these formulas
specify $G,k$. The positive definiteness of $\gamma$ follows already
from its convex interpolation. The estimates below give $d,u>0$
uniformly for large $R$. Also $v_0/r=1+O(r^{-3})$, so this interpolation
keeps $v>0$: no passage through $v=0$ occurs in the collar.
Where $w=e=0$, we have $c_R=0$, $u=u_0$, and the original data are
recovered exactly. The flat endpoints of the cutoffs then give a smooth
join to the unchanged interior. Smoothness here uses the smoothness
of the input; the uniform estimates will use only its stated finite
derivative bounds.

\medskip
\noindent\emph{Weighted estimates at the start of the modification.}
Write $D=r\partial_r$.  For an angular tensor $T$, the notation
$T=O^j_\omega(wr^{-a})$ will mean
\[
 \sum_{i=0}^j|\nabla_\sigma^iT|_\sigma\leq C_*w(r)r^{-a};
\]
this notation makes no assertion about radial derivatives.  In this
proof $\delta$ denotes the new quantity minus the original one; for
$L_b$ the original quantity is $L_0$.

The separated supports of the cutoffs imply
\begin{equation}\label{eq:collar-cutoff-estimates}
 |D^j e|\leq C_jw\quad(j\geq0),\qquad
 c_R(r)\leq2w(r)R^{-1/2}\leq2\sqrt3\,w(r)r^{-1/2}
 \quad(R\leq r\leq3R).
\end{equation}
For the first assertion, $e$ and all its derivatives vanish below $2R$,
whereas $w=1$ above $2R$.  For the second, use monotonicity to bound
$w(\rho)$ by $w(r)$ inside the integral defining $c_R$.  No bound on
$w'/w$ is used.

The original estimates \eqref{eq:original-asymptotics}, the interpolation
formulas, and \eqref{eq:collar-cutoff-estimates} give the following table.
The two columns mean a bound for the listed expression and for its
difference from the corresponding original expression, respectively.
\begin{equation}\label{eq:collar-weighted-table}
\begin{array}{c|c|c}
 \text{expression}&\text{bound}&\text{difference bound}\\ \hline
 r^{-2}\gamma-\sigma,\ d-1
   &O^2_\omega(r^{-3})&O^2_\omega(wr^{-3})\\
 u/\sqrt{1+r^2}-1,\ U/\sqrt{1+r^2}-1
   &O^2_\omega(r^{-3})&O^2_\omega(wr^{-3})\\
 v/r-1,\ L_b-1,\ \eta/r,\ \tau/r^2
   &O^2_\omega(r^{-3})&O^2_\omega(wr^{-3})\\
 s&O^2_\omega(r^{-5})&O^2_\omega(wr^{-5}).
\end{array}
\end{equation}
Here and below an omitted factor $w$ means the corresponding unweighted
estimate.  We justify the derivative assertions needed from this table.
The interpolation of $\gamma/r^2$ and $s$ has the same weighted bounds
through three total angular and scaled radial derivatives, by the three
derivatives available for $g$.  This gives the two angular derivatives of
$d$, because
\[
 d=\frac r4\tr_\gamma(\partial_r\gamma-\mathcal L_s\gamma)
\]
uses one such derivative.  For $v/r-1$ and $\eta/r$ the weighted bounds
also hold after one application of $D$, including the lower angular
orders allowed by the two derivatives of $K$.  The same holds for
$\gamma/r^2-\sigma$ and $s$.  These are the only radial derivative bounds
from the table that will be needed.

For completeness, the mass interpolation has the more informative
bounds
\begin{equation}\label{eq:collar-mass-difference}
 \delta F=O^2_\omega(wr^{-1/2}),\qquad
 \partial_r\delta F=wr^{-3/2}+O(wr^{-2}).
\end{equation}
Indeed,
\[
 \delta F=e(f_*-F_0)+c_R,\qquad
 \partial_r\delta F=e'(f_*-F_0)-e\partial_rF_0+wr^{-3/2},
\]
and $F_0-f_*=O_2(r^{-1})$.  Also $F$ and its first two angular
derivatives are uniformly bounded.  The bounds for $u$ in
\eqref{eq:collar-weighted-table} now follow algebraically from
\[
 \frac{u^2}{1+r^2}-1
   =\frac{v^2-r^2-2F/r}{1+r^2},\qquad
 \delta(u^2)=(v+v_0)\delta v-2\delta F/r.
\]
Thus $u/\sqrt{1+r^2}$ is positive and close to one.  The already proved
bound $d=1+O(r^{-3})$ gives $d>0$ and the same estimates for $U=u/d$.
Finally, inverse metrics and their angular derivatives preserve the
weighted bounds, so the trace removal in \eqref{eq:collar-data} gives the
claimed bounds for $\tau$.  None of these steps requires a weighted
radial derivative estimate for $u$, $U$, or $d$.

These estimates also give $|G-h|_h\leq C_*r^{-3}$.  For the mixed term in
\eqref{eq:foliation}, use $\gamma/r^2=O(1)$ and $s=O(r^{-5})$; for the
radial term use $U/\sqrt{1+r^2}=1+O(r^{-3})$.
To obtain the comparison with the original metric, choose $C_{\mathrm{met},0},C_{\mathrm{met},1}$
so that $g\leq(1+C_{\mathrm{met},0}R^{-3})h$ and
$G\geq(1-C_{\mathrm{met},1}R^{-3})h$ throughout the collar. For sufficiently large
$R$,
\[
 G\geq\frac{1-C_{\mathrm{met},1}R^{-3}}{1+C_{\mathrm{met},0}R^{-3}}g
   \geq\bigl(1-(C_{\mathrm{met},0}+C_{\mathrm{met},1})R^{-3}\bigr)g.
\]
Inside $r=R$ there is equality $G=g$. This proves the metric assertion
on the whole constructed region. We must still prove DEC there.

\medskip
\noindent\emph{The momentum increments.}
We compare the constraint quantities of Lemma~\ref{lem:constraints}
before and after the modification.  Their radial components satisfy
\begin{equation}\label{eq:collar-momentum-estimates}
 Q_0=O(r^{-3}),\qquad
 \delta Q=H\frac{w}{r^2}+O(wr^{-3}),\qquad H=2+O(r^{-2}),
\end{equation}
and the tangential components satisfy
\begin{equation}\label{eq:collar-transverse-estimates}
 |Z_0|_{\gamma_0}=O(r^{-3}),\qquad
 |Z-Z_0|_\gamma=O(wr^{-3}).
\end{equation}
Here the subtraction in the last formula is of angular covectors in the
fixed end coordinates, so is well defined before comparing their norms.
We give the termwise estimates to keep track of the factor $w$.

Taking the trace-free part of
$\chi=\tfrac U2(\partial_r\gamma-\mathcal L_s\gamma)$ gives
\[
 \hat\chi/r^2=O(r^{-3}),\qquad
 \delta(\hat\chi/r^2)=O(wr^{-3}).
\]
In fact the contribution from differentiating the factor $r^2$ in
$\gamma=r^2(r^{-2}\gamma)$ is $(U/r)\gamma$, which has no trace-free
part.  The acceleration covector
$B_N^\flat=\mathrm d_{\Sph}\log U$ is $O(r^{-3})$ with difference
$O(wr^{-3})$.  For the trace term put $t=v/r$ and $t_0=v_0/r$. The interpolation
is $t=(1-e)t_0+e$, so
\[
 \delta t=-e(t_0-1),\qquad
 \partial_r\delta t=-e'(t_0-1)-e\partial_rt_0=O(wr^{-4}).
\]
Here $t_0-1=O_2(r^{-3})$ and $|re'|\leq Cw$. Together with
$p=2t$, the angular bounds for $\delta t$, and the estimates for $U,s$,
this gives
\[
 Np=O(r^{-3}),\qquad \delta(Np)=O(wr^{-3}).
\]
In the formula for $Q$ in \eqref{eq:constraints}, its first term has
increment
\[
 \delta\bigl(H(L-p/2)\bigr)
   =H\frac{w}{r^2}+O(wr^{-3}).
\]
The shear contraction is quadratic in tensors of norm $O(r^{-3})$,
while $\diver_\gamma\eta=O(r^{-4})$ with difference $O(wr^{-4})$.
The acceleration term is smaller.  This proves
\eqref{eq:collar-momentum-estimates}.

For $Z$, work first in angular covector components.  The expression
\[
 U(\partial_r\eta-\mathcal L_s\eta)+H\eta
\]
is $O(r^{-2})$ with difference $O(wr^{-2})$: here
$\eta=O(r^{-2})$, $\partial_r\eta=O(r^{-3})$, and $U=O(r)$.
Both $\diver_\gamma\tau$ and
$\mathrm d_{\Sph}(L+p/2)$ are $O(r^{-3})$, with differences $O(wr^{-3})$.
In the latter assertion the added radial entry causes no difficulty,
because $\mathrm d_{\Sph}(w/r^2)=0$.  Finally,
$(L\gamma-k_T)(B_N,\cdot)$ obeys at least these bounds by
\eqref{eq:collar-weighted-table}.  Passing from angular components to a
$\gamma$-covector norm introduces a factor $O(r^{-1})$.  These facts
prove \eqref{eq:collar-transverse-estimates}.  Notice in particular that
no radial derivative of $L$ occurs in either momentum formula.

\medskip
\noindent\emph{The small null increment.}
Put $a=v/u$ and $a_0=v_0/u_0$.  The mass identity
\eqref{eq:mass-identity}, in which $L$ has canceled, gives
\begin{equation}\label{eq:collar-null-increment}
 \delta(C-aQ)=2wr^{-7/2}+O(wr^{-4}).
\end{equation}
In this expression $\delta(C-aQ)$ means
$(C-aQ)-(C_0-a_0Q_0)$.  To verify it, first note the exact identity
\[
 \frac{NF}{u}=\frac1d(\partial_rF-s^A\partial_AF).
\]
Using \eqref{eq:collar-mass-difference}, $\partial_rF_0=O(r^{-2})$,
$s=O(r^{-5})$, and the weighted bounds for $d,s$ therefore gives
\begin{equation}\label{eq:collar-mass-derivative}
 \delta\left(\frac{2NF}{ur^2}\right)
   =2wr^{-7/2}+O(wr^{-4}).
\end{equation}
The remaining terms in \eqref{eq:mass-identity} have the following
difference estimates:
\[
\begin{array}{c|c}
 \text{term}&\text{difference bound}\\ \hline
 R_\gamma/2-r^{-2}&O(wr^{-5})\\
 (d^{-1}-1)(2-6F/r)/r^2&O(wr^{-5})\\
 U\Delta_\gamma(U^{-1})&O(wr^{-5})\\
 -\tfrac12(|\hat\chi|^2+|\tau|^2)+a\hat\chi:\tau-|\eta|^2
   &O(wr^{-6})\\
 -a\bigl(\diver_\gamma\eta-2\eta(B_N)\bigr)&O(wr^{-4}).
\end{array}
\]
For the curvature term write
$R_\gamma=r^{-2}R_{r^{-2}\gamma}$ and use the two angular derivatives
of the normalized metric.  The second row uses $d-1=O(r^{-3})$ and
the boundedness of $F$.  For the third row put
$q=U/\sqrt{1+r^2}$; angular differentiation does not act on the radial
factor, so $U\Delta_\gamma(U^{-1})=q\Delta_\gamma(q^{-1})$, with
$q-1=O^2_\omega(r^{-3})$ and
$\delta q=O^2_\omega(wr^{-3})$.  The quadratic row follows from the
$O(r^{-3})$ sphere norms of $\hat\chi,\tau,\eta$ and their weighted
differences.  The last row uses the divergence estimate above.  Changes
of inverse metrics in these contractions have relative size
$O(wr^{-3})$ and respect all the displayed bounds.  This proves
\eqref{eq:collar-null-increment}.

In particular, neither the scalar-constraint calculation nor the
momentum estimates require $w'$: the mass term uses exactly
$c_R'=wr^{-3/2}$, angular derivatives ignore $w/r^2$, and the radial
derivative of $L$ has canceled from the constraints.  This is why the
bounds remain uniform at a flat start of the cutoff.

\medskip
\noindent\emph{The dominant energy cone.}
The reconstruction gives
\[
 a-1=O(r^{-2}),\qquad \delta a=O(wr^{-3}).
\]
Consequently
\[
 \delta(C-Q)
 =\delta(C-aQ)+(a-1)\delta Q+(a-a_0)Q_0
 =2wr^{-7/2}+O(wr^{-4}).
\]
For clarity, fix constants $A_*,B_*$, independent of $R$, such that
the preceding estimates give
\[
 \left|\delta(C-Q)-2wr^{-7/2}\right|\leq A_*wr^{-4},
 \qquad
 \left|\delta Q-2wr^{-2}\right|\leq B_*wr^{-3}.
\]
Taking $R\geq\max\{1,A_*^2,B_*\}$ yields
$\delta(C-Q)\geq wr^{-7/2}$ and $\delta Q\geq wr^{-2}$.
In particular, uniformly throughout the collar,
\begin{equation}\label{eq:collar-cone-margins}
 \delta C-\delta Q\geq wr^{-7/2}\geq0,
 \qquad
 \delta C+\delta Q\geq wr^{-2}\geq0.
\end{equation}

It remains to compare the tangential momenta in a common norm.  At each
point choose a linear isometry
\[
 I:(T^*\Sph,\gamma_0^{-1})\longrightarrow
       (T^*\Sph,\gamma^{-1})
\]
as follows. On covectors let $A$ be the positive self-adjoint operator,
relative to $\gamma_0^{-1}$, determined by
\[
 \gamma^{-1}(\alpha,\beta)=\gamma_0^{-1}(A\alpha,\beta).
\]
Set $I=A^{-1/2}$. Then
$\gamma^{-1}(I\alpha,I\beta)=\gamma_0^{-1}(\alpha,\beta)$.
Since $A-\mathrm{Id}=O(wr^{-3})$, spectral calculus on a fixed interval
around $1$ gives $I-\mathrm{Id}=O(wr^{-3})$ in the old covector norm.
This comparison is pointwise and needs no derivatives of $I$. By
\eqref{eq:collar-transverse-estimates} the covector
$W=Z-IZ_0$ satisfies
\[
 |W|_\gamma\leq C_1wr^{-3}.
\]
The product of the two lower bounds in
\eqref{eq:collar-cone-margins} is $w^2r^{-11/2}$. Increasing $R$ also
to satisfy $R\geq C_1^4$ gives
\[
 (\delta C)^2-(\delta Q)^2
 \geq w^2r^{-11/2}\geq C_1^2w^2r^{-6}\geq|W|_\gamma^2.
\]
The nonnegativity in \eqref{eq:collar-cone-margins} gives $\delta C\geq0$,
so the last inequality yields
\[
 \delta C\geq\sqrt{(\delta Q)^2+|W|_\gamma^2}.
\]
This argument also includes $w=0$, where all increments vanish exactly.
The old dominant energy condition and the triangle inequality in
$\R\oplus T^*\Sph$ now give
\begin{align*}
 C=C_0+\delta C
 &\geq\sqrt{Q_0^2+|IZ_0|_\gamma^2}
       +\sqrt{(\delta Q)^2+|W|_\gamma^2}\\
 &\geq\sqrt{Q^2+|Z|_\gamma^2}.
\end{align*}
This proves DEC even where the original constraint vector lies on
the boundary of the future cone. It also explains why estimating only
the scalar constraint would have been insufficient: the angular
momentum is controlled by the product of the two null margins.

Finally, all derivatives of $1-e$ vanish at $3R$, and $w=1$ there.
Thus the nonround leaf and mixed fields disappear to every order,
$v$ agrees with $r$ to every order, and $d$ agrees with $1$ to every
order. Equations~\eqref{eq:collar-data}--\eqref{eq:collar-reconstruction}
therefore give exactly the stated jets. The padding continues smoothly
to all $r\geq3R$, with $c_R'=r^{-3/2}$ and
$0\leq c_R\leq2R^{-1/2}$. These are the round data to which the
outer bending construction applies.
\end{proof}
